% page 339 du fac-simile (image p-338 du scan)
% bloc 165.1 x 242.0 mm ; marge gauche 36.9 mm
\pgc{15.240mm}{0.000mm}{
\l{\cel{53}331}
\l{}
\l{\sou{lasho}. I. (\sou{kirurg}.)Plako de ligno, de kartono, quan onu enpozas}
\l{\cel{3}alonge membro ruptita, por mantenar la osti en plaso fixa.}
\l{\cel{3}II. (\sou{tekn}.)Plako de fero qua kunjuntas du reli. - D.}
\l{}
\l{\sou{laskaro.}\cel{1}(\sou{navig.)}\cel{2}Membro Indiana di la +kruo di navo. - DEFIS.}
\l{}
\l{\sou{lasta}. I. Qua esas pos omni cetera, (l) pri tempo, (2) pri spaco,}
\l{\cel{3}(3) pri rango en serio. II. (\sou{+ultima)}\cel{1}: Pos qua ne plus existo}
\l{\cel{3}ceteri, (l) pri tempo, (2) pri spaco, (3) pri rango en serio.}
\l{\cel{3}DE.}
\l{}
\l{\sou{lato}.\cel{2}Peco de ligno fendita segun la longeso di la fibri, longa,}
\l{\cel{3}dina e plata, uzata por facar parieti, plafoni, e c. e qua}
\l{\cel{3}suportas la ardezo-plaki o la teguli di tekto. - DEFS.}
\l{}
\l{\sou{latanio.}\cel{2}(\sou{bot.}) Palmiero kun larja folii dispozita abanike, quan}
\l{\cel{3}onu vidas en Madagaskar ed en la insuli ri Indonezia. L.\cel{1}\sou{lata}-}
\l{\cel{3}\sou{nia}. - EFI.}
\l{}
\l{\sou{latenta}. Qua ne su manifestas adextere. -\cel{1}\sou{Kaloro latenta}\cel{1}: kaloro}
\l{\cel{3}quan absorbas korpo por liquideskar o vaporeskar, qua, pro ne}
\l{\cel{3}manifestar su adextere, ne efikas a la termometro. - DEFIS.}
\l{}
\l{\sou{latero}. I. (\sou{anat.})\cel{2}La tota parto, dextra o sinistra, di la}
\l{\cel{3}korpo homala. - II. Facio, parto qua esas dextre o sinistre}
\l{\cel{3}che objekto. Singla ek la facii interopozanta di objekto. -}
\l{\cel{3}III. (\sou{geom.}) Linei qui limitizas surfaco. - EFIS.}
\l{}
\l{\sou{latexo}. (\sou{bot.)}\cel{2}Suko propra di ula planti (euforbi,figieri, papa-}
\l{\cel{3}veri, latugi, e c.) e qua cirkulas en vaskuli partikulara. -}
\l{\cel{3}DEF.}
\l{}
\l{l\cel{1}\sou{a t i r}\cel{1}o. (\sou{bot.)}\cel{2}Planto klimera, di qua la flori esas odoroza}
\l{\cel{3}- DEF.}
\l{}
\l{\sou{laticifera}. (\sou{bot.)}\cel{3}Qua produktas latexo. - DEF.}
\l{}
\l{\sou{latitudo}. I. (\sou{kosmogr.)}\cel{2}Disto de loko ad equatoro, mezurata per}
\l{\cel{3}arko di la meridiano terala. - II. Regiono submisata a ta o}
\l{\cel{3}ca rejimo atmosferala, segun ke lu esas plu o min for equatoro}
\l{\cel{3}-EFIRS.}
\l{}
\l{\sou{latrino}\cel{1}lnoarezervita por kakar od urinifar. - DEFIS.}
\l{}
\l{latugo. (\sou{bot.)}\cel{3}Planto leguma, qua kontenas suko laktatra, ek la}
\l{\cel{3}tribuo \textquotedbl{}chikorei\textquotedbl{}. - L.\cel{1}\sou{lactuca}. - IS.}
\l{}
\l{latuno. Metalo ek aloyuro de kupro e de zinko. -\cel{2}DFRS.}
\l{}
\l{laubo. Treliso-muri paralela, kun mi-cilindro trelisa kom tekto,,}
\l{\cel{3}kovrita per planti klimera tre folioza. - D.}
\l{}
\l{laudar. (\sou{trans.})\cel{2}Igar eminenta per la aprobo di lua meriti, di}
\l{\cel{3}lua hona qualesi. - EFIS.}
}
